\section{Introduction}
A serving scheduler can reduce average response time while making long
requests slower or lowering aggregate throughput. Which scheduler appears
best therefore depends on both the workload and the objective. There is a
more basic difficulty in simulation-based comparisons: a policy label can
select not only an ordering rule, but also a different prefill model,
memory constraint, or information interface. An observed rank difference
then combines scheduling decisions with differences in what the evaluator
allows those decisions to accomplish.

This problem arises in the historical LLM-Serving Scheduler Portability
(LSSP) corpus that motivates our study. Its implementations shared request
inputs, but did not uniformly enforce prefill work and KV capacity. Some
comparisons therefore concerned policy--model combinations rather than
scheduling alternatives within a common execution model. Repeating those
runs or reproducing their ranking arithmetic cannot resolve that ambiguity.
A useful comparison must first identify the execution contract: which
requests are visible, which actions are legal, how work advances, and when
resources are acquired and released.

We evaluate four adaptations under such a contract, \cm. It is a CPU
discrete-event model with common batch timing, token accounting,
information restrictions, and conservative full-context KV reservations.
The adaptations retain selected ideas from FCFS serving, chunked prefill,
decode-first scheduling, and SLAI-inspired prompt ordering
\citep{kwon2023vllm,agrawal2024sarathi,slai2025}. They have new implementation
identities and a shared action space. The purpose is to measure their
behavior under a controlled abstraction, not to reproduce the performance
of their published systems.

Workload-sensitive serving evaluations already establish that benchmark
choice matters \citep{papaioannou2024workload}; predictive serving simulators
address the cost of exploring configurations \citep{vidur2024,llmservingsim2024}.
Our contribution is a different empirical object: a complete comparison
of restricted adaptations across a fixed corpus, with absolute performance,
matched differences, ties, and practical reversals reported together.
This permits a reader to distinguish a useful performance difference from
a rank change caused by aggregation, duplicate metrics, or an inactive
scheduling distinction.

The completed comparison changes the emphasis of the historical narrative.
Across six source--capacity aggregates, SLO-gated improves mean latency but
reduces throughput; BurstGPT also has worse mean window-p95 latency.
Chunked-FCFS and decode-first are exactly equivalent in every matched
condition. Source-mean comparisons have no strict 10\% practical reversals,
whereas individual windows and different metrics do. These findings make
an unconditional ``best scheduler'' an inadequate summary, without implying
that rankings are unstable at every aggregation level.

We ask three questions:
\begin{description}[leftmargin=0pt,labelindent=0pt]
\item[RQ1: Performance and operating conditions.] How do the absolute
performances of four scheduler adaptations vary across workload windows,
inherited arrival/deadline settings, and KV reservation capacities?
\item[RQ2: Ranking stability and metric choice.] How do workload variation,
aggregation, metric selection, ties, and practical performance differences
affect scheduler rankings?
\item[RQ3: Interpretation and inference.] What do the historical comparison
audit and known-truth statistical calibration establish about the reliability
and limitations of scheduler comparisons?
\end{description}

The paper makes four specific contributions. First, it specifies an
executable comparison contract and identifies the departures of each
adaptation from its published lineage. Second, it reports the complete
5,760-cell evaluation, including absolute latency, throughput, goodput,
and complete request accounting. Third, it quantifies how window variation,
aggregation, ties, and metric choice change the interpretation of rankings,
including the null source-mean reversal result. Fourth, it separates
execution conformance from statistical calibration through a concise
historical audit and a completed synthetic diagnostic. These contributions
support finite-corpus conclusions under \cm; they make no new algorithmic,
statistical-test, or native-performance claim.
